I wanted to know why training-data filtering removes some fine-tuned
behaviors but barely touches others. In the comments on the
Lee-Rosser-Engels-Nanda post,
\href{https://www.lesswrong.com/posts/aTybJ6CPQrxEY8rE2/data-filtering-works-a-lot-worse-than-you-would-expect?commentId=jmZwaZePbQJrtWvHS}{draganover}
proposed that fine-tuning can strengthen a whole bundle of
persona-linked behaviors.
\href{https://www.lesswrong.com/posts/aTybJ6CPQrxEY8rE2/data-filtering-works-a-lot-worse-than-you-would-expect?commentId=cqeJvGFmpmD66Hrer}{Dohun
Lee} discussed this as a possible explanation. I wanted to test a
concrete version of it. This connects my data-attribution background
with my interest in safety-relevant personas: when can identifying
influential data help us change behavior?

I reconstructed the filtering setup with
\texttt{allenai/Olmo-3-1025-7B}, rank-64 LoRA, and 25,000
Dolci-Think-SFT examples. Across 14 filter-retrain experiments, I ranked
training examples, deleted selected examples, retrained from the same
base, and measured behavior. Before attribution, I fixed four directions
in the model's layer-22 activations as my representation of persona
variation.

The pre-registered order held: refusal was easiest to remove, followed
by feelings validation, then bold formatting, with both-sides framing
and ethical frameworks hardest. The last two were one unordered group.
All nine comparisons between the predicted groups held, but two results
were already known when I fixed the final set, and each retrain used one
seed. This supports that declared ordering; the numerical
persona-separability measure itself did not recover the same order.

The proposed mechanism failed its direct tests. Behaviors expected to
share a persona did not show the predicted extra alignment with the
persona directions or overlap in highly ranked documents (H1). I then
subtracted the persona component from the rule used to rank training
examples and repeated deletion and retraining (H3). I had required at
least +1.0 judge point of improvement in feelings validation, with an
interval excluding zero. Instead, the estimated gain was
\textbf{-0.430}: more validation remained than under ordinary filtering.
The 95\% interval {[}-0.975, +0.105{]} includes zero, so this failed to
improve filtering without proving it became worse. The both-sides
control's predicted lack of improvement held. These results constrain
the specific frozen linear representation I tested; other
representations or training-time mechanisms remain possible.

Next I'd hold prompts and training budgets fixed while comparing
deletion with neutral rewrites, alternate-teacher completions, and
anti-trait completions. With random deletion as a control, this could
distinguish effects of completion content from data removal and guide
tests of training dynamics. Document coherence remains future work.

This is one model and one seed per retrain, with one final judge after
cross-provider revalidation failed. The reconstructed recipe and
behavior labels have explicit limitations, including different base/SFT
generation conventions. I reused
\href{https://arxiv.org/abs/2606.19625}{Capability Provenance}'s attribution
infrastructure, \href{https://arxiv.org/abs/2608.00023}{Role Steering}'s
275-role inventory, and existing GPU access. Agents wrote much of the code
and initial prose. I drove hypothesis generation and selection,
inspected outputs and ratings, and reviewed HPC and persona code most
closely. I did not audit every automated loop. This summary uses my
notes with AI drafting assistance.
